\subsubsection{Perceptual Similarity vs Pixel Similarity}\label{sec:perceptual-similarity-vs-pixel-similarity}

When we say that two images are ``\emph{similar}'', we must ask: \emph{\textbf{similar in what sense?}}; \emph{\textbf{Similar according to whom?}} There are actually two completely different notions of similarity:
\begin{enumerate}
    \item \definition{Pixel Similarity}: two images are pixel-similar if corresponding pixels have similar values.
    
    \highspace
    A \textbf{computer} initially sees an image only as a vector, for example, for CIFAR-10, a vector $x \in \mathbb{R}^{3072}$. The \hl{simplest way to compare two images is to compare their pixels}:
    \begin{equation*}
        \text{Pixel Similarity}(x_1, x_2) = \left\|x_1 - x_2\right\|_{2}
    \end{equation*}
    Or another pixel-wise distance.
    \begin{itemize}
        \item If most pixels are identical, the distance is small.
        \item If many pixels differ, the distance is large.
    \end{itemize}
    The computer therefore concludes: ``\emph{these images are similar because their pixels are similar}''.
    
    
    \item \definition{Perceptual Similarity} (or \definition{Semantic Similarity}): ignores the exact pixel values and focuses on \textbf{what the image represents}.
    
    \highspace
    Humans work very differently. We \textbf{don't compare images pixel by pixel}. Instead, we \hl{compare objects, shapes and structures, as well as semantics.} To us, two different breeds of dog are clearly similar, even though they are made up of completely different pixels. Humans still perceive them as similar because they share the same \emph{concept} of ``dog''.
\end{enumerate}
\textcolor{Green3}{\faIcon{question-circle} \textbf{Why are these two notions of similarity so different?}} Let's revisit everything we have said about the difficulty of image classification. Suppose we have a dog image. Now create five new images by applying the following transformations to the original image:
\begin{enumerate}
    \item \textbf{Translation}: Move the dog to the left or right.
    \item \textbf{Rotation}: Rotate the dog by a few degrees.
    \item \textbf{Shadow}: Add a shadow to the dog.
    \item \textbf{Different Viewpoint}: Take a picture of the dog from a different angle.
    \item \textbf{Different Pose}: Take a picture of the dog in a different pose.
\end{enumerate}
Humans immediately say ``\emph{these are all pictures of the same dog}''. But \hl{mathematically, each transformed image has many different pixels. The Euclidean distance between pixel vectors may become very large.}

\highspace
Now consider another example. Take a picture of a ``Wolf'' and a picture of a ``Dog''. The pixel values may actually be closer than a transformed image of the same dog. Yet semantically, the first pair belongs to two different classes, while te second pair belongs to the same class. This shows that \textbf{pixel similarity and semantic similarity are fundamentally different concepts}.

\highspace
Every challenge we have discussed so far in this section is a manifestation of the same underlying problem: \textbf{images that humans perceive as similar may have very different pixel values}. Although our simple linear classifier was \hl{effective in representing pixel similarity}, it was \hl{completely ineffective in representing abstract concepts due to limitations in semantic similarity.}

\highspace
\hl{Again, this is why we need CNNs.} Instead of comparing single pixels, CNNs compare \textbf{features} of images. At the deepest layers, the representation is no longer based on raw RGB values. Instead, it encodes \textbf{semantic features}. Consequently, the similarity measured inside a CNN becomes much closer to \textbf{perceptual similarity} than to raw pixel similarity. This is why \hl{CNNs are dramatically better than linear classifiers.}